\documentclass[11pt]{article}
\usepackage[utf8]{inputenc}
\usepackage{amsmath,amssymb,booktabs,hyperref}

\title{Synthetic Lattice Sum Identity}
\date{}

\begin{document}
\maketitle

\noindent\textbf{Status:} Done with open items: 1 step checked numerically\\

\section{Assumptions and Sources}
\label{sec:setup}

User-supplied lattice conventions for the finite chain.

\noindent\textbf{e-1} [Given]: Let N >= 2 be an integer and let x > 0.

\begin{quote}
The discrete boundary sum equals 7/12 for every even chain length.
\end{quote}
\noindent --- \cite{voss2024lattice} (p. 2, line 4) (OCR)

\noindent\textbf{e-2} [Found]: Boundary sum lemma from the cited lattice note.

\section{Derivation}
\label{sec:derivation}

We expand the polynomial identity and test the tail bound. \textit{(model)}

% handle: poly_id@1
\begin{equation}
\label{eq:poly}
(x + 1)^2 - (x^2 + 2x + 1) = 0
\end{equation}

\noindent\textbf{e-3} [Verified]: Symbolic expansion of the square vanishes identically.

% handle: tail_err@1
\[
\left|\sum_{k=1}^{64} k^{-2} - \frac{\pi^2}{6}\right| < 0.016
\]

\noindent\textbf{e-4} [Numerically checked]: Partial sum matches the analytic limit at 12 points.

\section*{How each statement was checked}

\begin{tabular}{p{0.18\linewidth}p{0.28\linewidth}p{0.46\linewidth}}
\toprule
Step / Claim & How checked & Rests on \\
\midrule
e-1 & Given & --- \\
e-2 & Found & --- (OCR source) \\
e-3 & Verified & e-1 \\
e-4 & Numerically checked & e-1, e-3 \\
\bottomrule
\end{tabular}

\bibliographystyle{plain}
\bibliography{writeup_sample}

\end{document}
